% VI47-DETAILED-HINTS
% VI47-DETAILED-HINT-01
\paragraph{Exercise VI.47.1.}
\begin{hint}
% PEDAGOGY-ENRICHED-VI
Use \(s_1-s_0\) on the overlap.
\end{hint}

% VI47-DETAILED-HINT-02
\paragraph{Exercise VI.47.2.}
\begin{hint}
% PEDAGOGY-ENRICHED-VI
Expand the three pairwise differences.
\end{hint}

% VI47-DETAILED-HINT-03
\paragraph{Exercise VI.47.3.}
\begin{hint}
% PEDAGOGY-ENRICHED-VI
Use sheaf gluing.
\end{hint}

% VI47-DETAILED-HINT-04
\paragraph{Exercise VI.47.4.}
\begin{hint}
% PEDAGOGY-ENRICHED-VI
There are no higher increasing tuples.
\end{hint}

% VI47-DETAILED-HINT-05
\paragraph{Exercise VI.47.5.}
\begin{hint}
% PEDAGOGY-ENRICHED-VI
Use \(c_{12}-c_{02}+c_{01}=0\).
\end{hint}

% VI47-DETAILED-HINT-06
\paragraph{Exercise VI.47.6.}
\begin{hint}
% PEDAGOGY-ENRICHED-VI
The overlap section group is zero.
\end{hint}

% VI47-DETAILED-HINT-07
\paragraph{Exercise VI.47.7.}
\begin{hint}
% PEDAGOGY-ENRICHED-VI
Track one global section under refinement.
\end{hint}

% VI47-DETAILED-HINT-08
\paragraph{Exercise VI.47.8.}
\begin{hint}
% PEDAGOGY-ENRICHED-VI
Use all nonempty intersections \(U_i\cap V_j\).
\end{hint}

% VI47-DETAILED-HINT-09
\paragraph{Exercise VI.47.9.}
\begin{hint}
% PEDAGOGY-ENRICHED-VI
Extend the overlap section to one open.
\end{hint}

% VI47-DETAILED-HINT-10
\paragraph{Exercise VI.47.10.}
\begin{hint}
% PEDAGOGY-ENRICHED-VI
Check all finite intersections.
\end{hint}

% VI47-DETAILED-HINT-11
\paragraph{Exercise VI.47.11.}
\begin{hint}
% PEDAGOGY-ENRICHED-VI
Both \(t\) and \(t^{-1}\) are regular.
\end{hint}

% VI47-DETAILED-HINT-12
\paragraph{Exercise VI.47.12.}
\begin{hint}
% PEDAGOGY-ENRICHED-VI
Use exponents \(0,1,2\).
\end{hint}

% VI47-DETAILED-HINT-13
\paragraph{Exercise VI.47.13.}
\begin{hint}
% PEDAGOGY-ENRICHED-VI
The ranges \(\ge0\) and \(\le-1\) cover all exponents.
\end{hint}

% VI47-DETAILED-HINT-14
\paragraph{Exercise VI.47.14.}
\begin{hint}
% PEDAGOGY-ENRICHED-VI
The missing exponents are \(-2,-1\).
\end{hint}

% VI47-DETAILED-HINT-15
\paragraph{Exercise VI.47.15.}
\begin{hint}
% PEDAGOGY-ENRICHED-VI
Intersect \(k[t]\) and \(t^{-2}k[t^{-1}]\).
\end{hint}

% VI47-DETAILED-HINT-16
\paragraph{Exercise VI.47.16.}
\begin{hint}
% PEDAGOGY-ENRICHED-VI
Insert \(n=0\).
\end{hint}

% VI47-DETAILED-HINT-17
\paragraph{Exercise VI.47.17.}
\begin{hint}
% PEDAGOGY-ENRICHED-VI
A 1-cocycle is transition data.
\end{hint}

% VI47-DETAILED-HINT-18
\paragraph{Exercise VI.47.18.}
\begin{hint}
% PEDAGOGY-ENRICHED-VI
Multiply \(t^n\) and \(t^m\).
\end{hint}

% VI47-DETAILED-HINT-19
\paragraph{Exercise VI.47.19.}
\begin{hint}
% PEDAGOGY-ENRICHED-VI
Leray is sufficient, not necessary.
\end{hint}

% VI47-DETAILED-HINT-20
\paragraph{Exercise VI.47.20.}
\begin{hint}
% PEDAGOGY-ENRICHED-VI
Use the pairwise-intersection cover.
\end{hint}

% VI47-DETAILED-HINT-21
\paragraph{Exercise VI.47.21.}
\begin{hint}
% PEDAGOGY-ENRICHED-VI
Missing exponents run from \(n+1\) to \(-1\).
\end{hint}

% VI47-DETAILED-HINT-22
\paragraph{Exercise VI.47.22.}
\begin{hint}
% PEDAGOGY-ENRICHED-VI
Split into the three integer ranges.
\end{hint}

% VI47-DETAILED-HINT-23
\paragraph{Exercise VI.47.23.}
\begin{hint}
% PEDAGOGY-ENRICHED-VI
Recall the Cech/sheaf comparison warning.
\end{hint}

% VI47-DETAILED-HINT-24
\paragraph{Exercise VI.47.24.}
\begin{hint}
% PEDAGOGY-ENRICHED-VI
Subtract local lifts on overlaps.
\end{hint}
